\documentclass[11pt,a4paper]{article}
\usepackage[utf8]{inputenc}
\usepackage[T1]{fontenc}
\usepackage[italian]{babel}
\usepackage{lmodern,amsmath,amssymb,microtype}
\usepackage[a4paper,top=1.55cm,bottom=1.55cm,left=1.65cm,right=1.65cm]{geometry}
\usepackage{xcolor,enumitem,fancyhdr,needspace}
\usepackage[hidelinks]{hyperref}
\definecolor{blu}{RGB}{30,62,92}
\setlength{\parindent}{0pt}
\setlength{\parskip}{3pt}
\setlength{\headheight}{14pt}
\setlength{\abovedisplayskip}{5pt}
\setlength{\belowdisplayskip}{5pt}
\setlength{\abovedisplayshortskip}{4pt}
\setlength{\belowdisplayshortskip}{4pt}
\setlength{\jot}{2.5pt}
\everymath{\displaystyle}
\everydisplay{\displaystyle}
\setlist[enumerate]{leftmargin=*,itemsep=2pt,topsep=3pt}
\pagestyle{fancy}\fancyhf{}
\fancyhead[L]{\small Classe 2 BS | Equazioni della fila A}
\fancyhead[R]{\small Anno scolastico 2026/27}
\fancyfoot[C]{\small\thepage}
\renewcommand{\headrulewidth}{0.3pt}
\newcommand{\R}{\mathbb R}
\newcommand{\tit}[1]{{\Large\bfseries\color{blu}#1}\par\smallskip}
\newcommand{\sott}[1]{\par\medskip{\large\bfseries\color{blu}#1}\par\smallskip}
\newcommand{\nota}[1]{\par\smallskip\noindent\colorbox{blu!6}{\parbox{\dimexpr\linewidth-2\fboxsep\relax}{#1}}\par\smallskip}
\begin{document}
\tit{Trova i tre errori}
\textbf{Nome e cognome:} \hrulefill\quad\textbf{Data:} \rule{2.6cm}{0.4pt}\par
In ogni svolgimento ci sono \textbf{esattamente tre errori}, distribuiti tra le righe.
Alcuni passaggi sono corretti. Lavora sul quaderno, nell’ordine seguente:
\begin{enumerate}
\item Controlla nell’equazione originale la conclusione proposta sull’insieme $S$.
      Se un valore funziona, chiediti anche se l’insieme contiene tutte le soluzioni.
\item Per ciascuno dei tre errori indica \textbf{tra quali righe nasce}, spiega la regola
      violata e scrivi il passaggio corretto.
\item Risolvi di nuovo l’equazione originale e scrivi l’insieme completo delle soluzioni.
\end{enumerate}
\textbf{Attenzione:} un dato sbagliato riportato coerentemente nelle righe successive è una
conseguenza del primo errore, non un nuovo errore. Una conclusione corretta non garantisce
che lo svolgimento sia corretto.
\sott{Esercizio 2a}
\textbf{Riga 0: equazione originale.}
\[
\frac{2x-1}{\displaystyle\frac12+\frac13}-\frac{x+2}{\displaystyle\frac34-\frac14}=\frac{x-4}{\displaystyle1-\frac16}
\]
\textbf{Svolgimento da esaminare.}
\[\begin{aligned}
\text{(1)}\quad & \frac{2x-1}{\displaystyle\frac25}-\frac{x+2}{\displaystyle\frac12}=\frac{x-4}{\displaystyle\frac56}\\
\text{(2)}\quad & \frac52(2x-1)-2(x+2)=\frac65(x-4)\\
\text{(3)}\quad & 25(2x-1)-20(x+2)=12(x-4)\\
\text{(4)}\quad & 50x-25-20x+40=12x-48\\
\text{(5)}\quad & 30x+15=12x-48\\
\text{(6)}\quad & 18x+15=-48\\
\text{(7)}\quad & 18x=-48+15\\
\text{(8)}\quad & 18x=-33\\
\text{(9)}\quad & x=-\frac{11}{6}\\
\text{(10)}\quad & \text{Determinata:}\quad S=\left\{-\frac{11}{6}\right\}.
\end{aligned}\]
\sott{Esercizio 2b}
\textbf{Riga 0: equazione originale.}
\[
(2x-1)^2-(x+2)(x-2)=3x(x-2)+2x+5
\]
\textbf{Svolgimento da esaminare.}
\[\begin{aligned}
\text{(1)}\quad & 4x^2+1-(x^2-4)=3x(x-2)+2x+5\\
\text{(2)}\quad & 4x^2+1-(x^2-4)=3x^2-6x+2x+5\\
\text{(3)}\quad & 4x^2+1-x^2-4=3x^2-6x+2x+5\\
\text{(4)}\quad & 3x^2-3=3x^2-4x+5\\
\text{(5)}\quad & -3=-4x+5\\
\text{(6)}\quad & -8=-4x\\
\text{(7)}\quad & x=\frac{-8}{4}\\
\text{(8)}\quad & x=-2\\
\text{(9)}\quad & \text{Determinata:}\quad S=\{-2\}.
\end{aligned}\]
\clearpage\tit{Trova i tre errori | continua}
\sott{Esercizio 2c}
\textbf{Riga 0: equazione originale.}
\[
\frac{(3x-1)^2-(3x+1)(3x-1)}{2}=-3(x+1)+\frac52
\]
\textbf{Svolgimento da esaminare.}
\[\begin{aligned}
\text{(1)}\quad & \frac{9x^2-6x+1-(9x^2-1)}{2}=-3(x+1)+\frac52\\
\text{(2)}\quad & \frac{9x^2-6x+1-9x^2-1}{2}=-3(x+1)+\frac52\\
\text{(3)}\quad & \frac{-6x}{2}=-3(x+1)+\frac52\\
\text{(4)}\quad & -3x=-3x-3+\frac52\\
\text{(5)}\quad & -6x=-6x-3+5\\
\text{(6)}\quad & -6x+6x=-3+5\\
\text{(7)}\quad & 0=-3+5\\
\text{(8)}\quad & 0=8\\
\text{(9)}\quad & \text{Impossibile:}\quad S=\varnothing.
\end{aligned}\]
\sott{Esercizio 2d}
\textbf{Riga 0: equazione originale.}
\[
(x^2+x-2)^2+(2x-1)(2x+1)=x^2(x+1)^2+4x+1
\]
\textbf{Svolgimento da esaminare.}
\[\begin{aligned}
\text{(1)}\quad & x^4+2x^3-3x^2+4+(2x-1)(2x+1)=x^2(x+1)^2+4x+1\\
\text{(2)}\quad & x^4+2x^3-3x^2+4+(2x-1)(2x+1)=x^2(x^2+2x+1)+4x+1\\
\text{(3)}\quad & x^4+2x^3-3x^2+4+(2x-1)(2x+1)=x^4+2x^3+x^2+4x+1\\
\text{(4)}\quad & x^4+2x^3-3x^2+4+4x^2+1=x^4+2x^3+x^2+4x+1\\
\text{(5)}\quad & x^4+2x^3+x^2+5=x^4+2x^3+x^2+4x+1\\
\text{(6)}\quad & 5=4x+1\\
\text{(7)}\quad & 4=4x\\
\text{(8)}\quad & x=4\cdot4\\
\text{(9)}\quad & x=16\\
\text{(10)}\quad & \text{Determinata:}\quad S=\{16\}.
\end{aligned}\]
\clearpage\tit{Trova i tre errori | continua}
\sott{Esercizio 4a}
\textbf{Riga 0: equazione originale.}
\[
x^2-7x+6=0
\]
\textbf{Svolgimento da esaminare.}
\[\begin{aligned}
\text{(1)}\quad & (x+7)(x-1)=0\\
\text{(2)}\quad & x+7=0\quad\text{oppure}\quad x-1=0\\
\text{(3)}\quad & x=0-7\quad\text{oppure}\quad x=0-1\\
\text{(4)}\quad & x=-7\quad\text{oppure}\quad x=-1\\
\text{(5)}\quad & S=\{-1\}.
\end{aligned}\]
\sott{Esercizio 4c}
\textbf{Riga 0: equazione originale.}
\[
x^5-2x^4-18x^3+36x^2+81x-162=0
\]
\textbf{Svolgimento da esaminare.}
\[\begin{aligned}
\text{(1)}\quad & x^4(x-2)-18x^2(x-2)+81(x-2)=0\\
\text{(2)}\quad & x(x-2)(x^4-18x^2+81)=0\\
\text{(3)}\quad & x(x-2)(x^2-9)^2=0\\
\text{(4)}\quad & x(x-2)(x-9)^2(x+9)^2=0\\
\text{(5)}\quad & x(x-2)[(x-9)(x+9)]^2=0\\
\text{(6)}\quad & \begin{cases}x=0\\x-2=0\\x-9=0\\x+9=0\end{cases}\\
\text{(7)}\quad & x=0\quad\text{e}\quad x=2\quad\text{e}\quad x=9\quad\text{e}\quad x=-9\\
\text{(8)}\quad & \text{Impossibile:}\quad S=\varnothing.
\end{aligned}\]
\end{document}